% Generated draft: requires mathematical and source review.
\documentclass[11pt]{article}
\usepackage[margin=25mm]{geometry}
\usepackage{fontspec}
\setmainfont{Latin Modern Roman}
\usepackage{amsmath,amssymb,mathtools,unicode-math}
\setmathfont{Latin Modern Math}
\usepackage{enumitem,needspace,xurl,hyperref}
\setlength{\parindent}{0pt}
\setlength{\parskip}{5pt}
\setlength{\emergencystretch}{3em}
\newcommand{\N}{\mathbb{N}}
\newcommand{\Z}{\mathbb{Z}}
\newcommand{\Q}{\mathbb{Q}}
\newcommand{\R}{\mathbb{R}}
\newcommand{\C}{\mathbb{C}}
\newcommand{\F}{\mathbb{F}}
\newcommand{\A}{\mathbb{A}}
\newcommand{\PP}{\mathbb{P}}
\newcommand{\E}{\mathbb{E}}
\newcommand{\eps}{\varepsilon}
\newcommand{\dd}{\,\mathrm{d}}
\begin{document}
\section*{UnsolvedMath ID: 4400003}
\textbf{Source problem code:} catalogue entry
\subsection{Definitions and mathematical statement}
Let $d\geq1$ be an integer, $A=\{0,1,\ldots,d-1\}$, and $X=A^{\Z}$ with its product Borel $\sigma$-field. Use the left shift $(\sigma x)_j=x_{j+1}$. For $n\geq0$, define count vectors $v_n,w_n:X\to\Z_{\geq0}^d$ by
\[
(v_n(x))_i=\#\{0\leq j\leq n:x_j=i\},\qquad
(w_n(x))_i=\#\{0\leq j\leq n:x_{-j}=i\},\quad i\in A.
\]
Thus each vector counts $n+1$ symbols and includes $x_0$. Write $\mathcal B(v_n,v_{n+1},\ldots)$ for the $\sigma$-field generated by those count maps. The source defines
\[
\mathcal F_n^+=\bigcap_{k\geq n}\mathcal B(v_n,v_{n+1},\ldots),\qquad
\mathcal F^+=\bigcap_{n\geq0}\mathcal F_n^+,
\]
with analogous minus definitions using $w_n$, and asks whether $\mathcal F^+=\mathcal F^-$.

In the printed definition of $\mathcal F_n^+$, $k$ does not occur in the intersected expression, so literally $\mathcal F_n^+=\mathcal B(v_n,v_{n+1},\ldots)$ and
\[
\mathcal F^+=\bigcap_{n\geq0}\sigma(v_m:m\geq n),\qquad
\mathcal F^-=\bigcap_{n\geq0}\sigma(w_m:m\geq n).
\]
The analogous intended expression with $v_k,v_{k+1},\ldots$ in the inner intersection would change the intermediate fields but gives the same final tail field.

No probability measure is supplied. Literal equality of Borel sub-$\sigma$-fields on the full shift is different from equality modulo null sets under a specified shift-invariant measure. The research summary suggests a measure-theoretic tail-field question, but does not specify such a measure, a class of measures, or completion conventions. These must be checked before the supplied open-status interpretation can be adopted.
\subsection{Short English statement}
Do the tail $\sigma$-fields generated by forward and backward cumulative symbol counts coincide? The printed inner intersection has an unused index, and the source does not say whether equality is literal or modulo null sets for a chosen invariant measure.
\subsection{Sources}
\begin{itemize}
\item Local UnsolvedMath registry, permanent ID 4400003, exact source code catalogue entry, ``Pingree open problems---Petersen tail-field problem 1''; proposed by Karl Petersen, 2010. Registry status open, classification OPEN-TRIAGE.
\item \emph{Open problems from the 3rd Pingree Workshop on Dynamical Systems} (2010), Karl Petersen, first Problem, \url{https://math.huji.ac.il/~mhochman/open-problems/pingree-open-problems.pdf}. Supplied pdf-text extraction accessed July 29, 2026; the document presents the item as open and no later authoritative resolution check is recorded. No PDF was accessed here.
\item The supplied research summary says no later authoritative resolution was located in its web search through 2026. It places the question in work on count-statistic tail fields and names Petersen and coauthors, including vague R\'enyi-type terminology, without precise references or hypotheses. These remarks do not supply an invariant measure or a theorem resolving the formulation.
\item Provenance flags status and rights NEEDS\_REVIEW: public access was verified, but redistribution permission was not established.
\end{itemize}
\end{document}
